\documentclass{article}
\usepackage{mathtools} 
\usepackage{fontspec}
\usepackage[UTF8]{ctex}
\usepackage{amsthm}
\usepackage{mdframed}
\usepackage{xcolor}
\usepackage{amssymb}
\usepackage{amsmath}


% 定义新的带灰色背景的说明环境 zremark
\newmdtheoremenv[
  backgroundcolor=gray!10,
  % 边框与背景一致，边框线会消失
  linecolor=gray!10
]{zremark}{说明}

% 通用矩阵命令: \flexmatrix{矩阵名}{元素符号}{行数}{列数}
\newcommand{\flexmatrix}[4]{
  \[
  #1 = \begin{pmatrix}
    #2_{11}     & #2_{12}     & \cdots & #2_{1#4}   \\
    #2_{21}     & #2_{22}     & \cdots & #2_{2#4}   \\
    \vdots      & \vdots      & \ddots & \vdots     \\
    #2_{#31}    & #2_{#32}    & \cdots & #2_{#3#4}
  \end{pmatrix}
  \]
}

% 简化版命令（默认矩阵名为A，元素符号为a）: \quickmatrix{行数}{列数}
\newcommand{\quickmatrix}[2]{\flexmatrix{A}{a}{#1}{#2}}

\begin{document}
\title{习题 7.1}
\author{张志聪}
\maketitle

\section*{3}

仿照定理7.1的证明。

(i) 存在性；

$\{a_n\}$为递增数列，且$a_n < b_1, n = 1, 2, \cdots$，
依单调有界定理，$\{a_n\}$有极限$\xi$，且有
\begin{align*}
  a_n \leq \xi, n = 1, 2, \cdots
\end{align*}
同理，递减有界数列$\{b_n\}$也有界限，并按照$\lim\limits_{n \to \infty} b_n - a_n = 0$，有
\begin{align*}
  \lim\limits_{n \to \infty} b_n = \lim\limits_{n \to \infty} a_n = \xi
\end{align*}
且
\begin{align*}
  \xi \leq b_n, n = 1, 2, \cdots
\end{align*}
综上，
\begin{align*}
  a_n \leq \xi \leq b_n, n = 1, 2, \cdots
\end{align*}
$\because a_n < a_{n + 1} \leq \xi \leq b_{n + 1} < b_n, n = 1, 2, \cdots$。 \\
$\therefore a_n < \xi < b_n, n = 1, 2, \cdots$。

(ii) 唯一性。

设数$\xi'$也满足
\begin{align*}
  a_n < \xi' < b_n, n = 1, 2, \cdots
\end{align*}
于是有
\begin{align*}
  a_n - b_n < \xi - \xi' < b_n - a_n, n = 1, 2, \cdots
\end{align*}
即
\begin{align*}
  |\xi - \xi'| < b_n - a_n, n = 1, 2, \cdots
\end{align*}
由题设$\lim\limits_{n \to \infty} b_n - a_n = 0$和保不等式性，有
\begin{align*}
  |\xi - \xi'| \leq 0
\end{align*}
故有$\xi = \xi'$。

\section*{5}

\begin{itemize}
  \item (1)

        对任意$x_0 \in (0, 1)$，
        需证明存在正整数$n_0$使得
        \begin{align*}
          \frac{1}{n_0 + 2} < x_0 < \frac{1}{n_0}
        \end{align*}
        即
        \begin{align*}
          n_0 < \frac{1}{x_0} < n_0 + 2 \ \ \ (I)
        \end{align*}
        取$n_0 = \left[\frac{1}{x_0}\right] - 1$，$(I)$式可成立。

  \item (2)

        (i)$\left(0, \frac{1}{2}\right)$不可以被有限覆盖；

        反证法，$\left(0, \frac{1}{2}\right)$可以被有限开覆盖，
        设这有限个开区间为
        \begin{align*}
          H' = \{\left(\frac{1}{n_i + 2}, \frac{1}{n_i}\right)| i = 1, 2, \cdots, k\}
        \end{align*}
        因为$H'$中的元素个数是有限的，于是可取
        \begin{align*}
          L = max\{n_1, n_2, \cdots, n_k\} \\
          R = min\{n_1, n_2, \cdots, n_k\}
        \end{align*}
        我们有
        \begin{align*}
          0 < \frac{1}{L + 2} < \frac{1}{R} < 1
        \end{align*}
        于是，可取$0 < \epsilon < \frac{1}{L + 2}$（稠密性保证这个$\epsilon$一定存在），
        显然，有$\epsilon \notin H'$，出现矛盾。

        (ii)$\left(\frac{1}{100}, 1\right)$可以被有限覆盖，具体被
        开集合
        \begin{align*}
          H = \{\left(\frac{1}{n + 2}, \frac{1}{n}\right)| n = 1, 2, \cdots, 98\}
        \end{align*}
        覆盖。

        任意$x_0 \in \left(\frac{1}{100}, 1\right)$， 即
        \begin{align*}
          \frac{1}{100} < x_0 < 1 \\
          1 < \frac{1}{x_0} < 100
        \end{align*}
        于是
        \begin{align*}
          1 \leq [\frac{1}{x_0}] \leq 99 \ \ \ (1)
        \end{align*}

        欲证明存在$n_0$使得
        \begin{align*}
          \frac{1}{n_0 + 2} < x_0 < \frac{1}{n_0}
        \end{align*}
        即
        \begin{align*}
          n_0 < \frac{1}{x_0} < n_0 + 2 \ \ \ (2)
        \end{align*}
        由(1)、(2)可知，这样的$n_0$一定存在。

\end{itemize}

\section*{7}

提示：反证法。

\section*{8}

因为$S$是有界无限点集，故存在$M > 0$，使得$S \subseteq [-M, M]$。
反证法，假设$[-M, M]$中没有$S$的聚点，由聚点定义可知，对任意$x \in [-M, M]$，
都存在一个邻域$U(x; \delta_x)$只含有
$S$中有限多个点。开区间
\begin{align*}
  H = \{(x; \delta_x)| x \in [-M, M]\}
\end{align*}
是闭区间$[-M, M]$的一个无限开覆盖，
由开覆盖定理可知，$H$中存在一个有限开区间覆盖$[-M, M]$，
又因为这有限个区间中每个都只有$S$中有限多个点，
于是可得$S$中的点是有限的，出现矛盾。

\section*{9}

只证明充分性，必要性通过极限定义即可证明。

设数列$\{a_n\}$满足柯西条件，易得$\{a_n\}$是有界数列，
于是由聚点定理可知，集合$\{a_n\}$至少有一个聚点。假设
聚点不唯一，即存在$\xi_1, \xi_2$都是聚点，不妨设$\xi_1 < \xi_2$，
令$\delta = \xi_2 - \xi_1$，由聚点的定义可知，$U(\xi_1; \frac{1}{3}\delta)$，
$U(\xi_1; \frac{1}{3}\delta)$都有$S$中无穷多个点。
由柯西条件可知，存在$N > 0$，使得只要$m, n > N$，有
\begin{align*}
  |a_m - a_n| < \frac{1}{3}\delta
\end{align*}
因为$U(\xi_1; \frac{1}{3}\delta)$，
$U(\xi_1; \frac{1}{3}\delta)$都有$S$中无穷多个点，
所以存在$a_{n_1} \in U(\xi_1; \frac{1}{3}\delta), a_{n_2} \in U(\xi_2; \frac{1}{3}\delta)$，
其中$n_1, n_2 > N$，于是我们有
\begin{align*}
  |a_{n_1} - a_{n_2}| > \frac{1}{3}\delta
\end{align*}
出现矛盾，所以$\{a_n\}$中的聚点是唯一的。

接下来，证明$\{a_n\}$的极限就是这个唯一聚点$\xi$。
对任意$\epsilon > 0$，因为$\xi$是聚点，所以存在$U(\xi; \epsilon)$中有$\{a_n\}$中无数个点，
且$(\xi; \epsilon)$外只有有限个点，
否则按照聚点定理会存在另一个聚点（显然，这个聚点不等于$\xi$），与聚点的唯一性产生矛盾。
取有限个点下标的最大值为$N$，于是$n > N$时，有
\begin{align*}
  |a_n - \xi| < \epsilon
\end{align*}
故，数列$\{a_n\}$收敛于$\xi$。

\section*{10}

\textbf{根的存在性定理：}

设$f$在区间$[a, b]$上连续，$f(a)f(b) < 0$，证明$f(x) = 0$在区间$[a, b]$有实数根。

\textbf{证明：}

反证法，假设$f$在区间$[a, b]$上无实数根，
由局部保号性可知，对任意$x \in [a, b]$，存在$U_x = U(x; \delta_x) \cap [a, b]$，使得
$x \in U_x$与$f(x)$保持同号，因此所有的$U(x; \delta_x)$构成$[a,b]$一个无限开覆盖，
由覆盖定理可知，存在一个有限开区间覆盖$[a, b]$，记为$H$。

$a \in H$，所以$a$属于某个开区间，设为$a \in (x_0 - \delta_0, x_0 + \delta_0)$，
右端点$x_0 + \delta_0$属于$H$中另一个开区间，否则这个开区间$(x_0 - \delta_0, x_0 + \delta_0)$将
包含$b$点，和保号性矛盾。设$x_0 + \delta_0 \in (x_1 - \delta_1, x_1 + \delta_1)$，
由保号性可知，任意$x \in (x_1 - \delta_1, x_1 + \delta_1)$，$f(x) < 0$，
依此类推，一直向后移动$n$次，直到开区间包含$b$点，记为$b \in (x_n - \delta_n, x_n + \delta_n)$，
由上述操作可知，任意$x \in (x_n - \delta_n, x_n + \delta_n)$，$f(x) < 0$，
这与题设$f(b) > 0$矛盾，假设不成立，命题得证。

\section*{11}

\textbf{一致连续性定理：}

若函数$f$在闭区间$[a, b]$上连续，则$f$在$[a, b]$上一致连续。

\textbf{证明：}

对任意$\epsilon > 0$，由$f$在闭区间$[a, b]$上连续可知，
任意$x \in [a, b]$，存在$U(x; \delta_x)$，只要$x_0 \in U(x; \delta_x)$，有
\begin{align*}
  |f(x) - f(x_0)| < \frac{\epsilon}{2}
\end{align*}
取
\begin{align*}
  H = \{(x; \frac{\delta_{x}}{2})| x \in [a, b]\}
\end{align*}
因此，$H$是闭区间$[a, b]$的一个无限开覆盖，
由有限覆盖定理可知，存在一个有限开区间覆盖$[a, b]$，记为
\begin{align*}
  S = \{(x_i; \frac{\delta_{x_i}}{2})| i = 1, 2, \cdots, n\}
\end{align*}
取$\delta = min\{\frac{\delta_{x_1}}{2}, \frac{\delta_{x_2}}{2}, \cdots, \frac{\delta_{x_n}}{2}\}$，
对任意$x', x'' \in [a, b], |x' - x''| < \delta$，
设$x' \in (x_i; \delta_{i})$，我们有
\begin{align*}
  x' - x_i < \frac{\delta_{x_i}}{2}
\end{align*}
又因为
\begin{align*}
  |x'' - x_i| & = |x'' - x' + x' - x_i|           \\
              & \leq |x'' - x'| + |x' - x_i|      \\
              & < \delta + \frac{\delta_{x_i}}{2} \\
              & \leq \delta_{x_i}
\end{align*}
所以
\begin{align*}
  |f(x') - f(x'')|
   & = |f(x') - f(x_i) + f(x_i) - f(x'')|      \\
   & \leq |f(x') - f(x_i)| + |f(x_i) - f(x'')| \\
   & < \frac{\epsilon}{2} + \frac{\epsilon}{2} = \epsilon
\end{align*}
命题得证。

\end{document}